%% 
%% Copyright 2007-2025 Elsevier Ltd
%% 
%% This file is part of the 'Elsarticle Bundle'.
%% ---------------------------------------------
%% 
%% It may be distributed under the conditions of the LaTeX Project Public
%% License, either version 1.3 of this license or (at your option) any
%% later version.  The latest version of this license is in
%%    http://www.latex-project.org/lppl.txt
%% and version 1.3 or later is part of all distributions of LaTeX
%% version 1999/12/01 or later.
%% 
%% The list of all files belonging to the 'Elsarticle Bundle' is
%% given in the file `manifest.txt'.
%% 
%% Template article for Elsevier's document class `elsarticle'
%% with harvard style bibliographic references

\documentclass[preprint,12pt]{elsarticle}
\usepackage{subcaption}
\usepackage{placeins}
\usepackage{hyperref}
%% Use the option review to obtain double line spacing
%% \documentclass[preprint,review,12pt]{elsarticle}

%% Use the options 1p,twocolumn; 3p; 3p,twocolumn; 5p; or 5p,twocolumn
%% for a journal layout:
%% \documentclass[final,1p,times]{elsarticle}
%% \documentclass[final,1p,times,twocolumn]{elsarticle}
%% \documentclass[final,3p,times]{elsarticle}
%% \documentclass[final,3p,times,twocolumn]{elsarticle}
%% \documentclass[final,5p,times]{elsarticle}
%% \documentclass[final,5p,times,twocolumn]{elsarticle}

%% For including figures, graphicx.sty has been loaded in
%% elsarticle.cls. If you prefer to use the old commands
%% please give \usepackage{epsfig}


%% The amssymb package provides various useful mathematical symbols
\usepackage{amssymb}
%% The amsmath package provides various useful equation environments.
\usepackage{amsmath}
\usepackage{xcolor}
\usepackage{booktabs}
\usepackage{tabularx}
\newcommand{\revblue}[1]{\textcolor{blue}{#1}}
%% The amsthm package provides extended theorem environments
%% \usepackage{amsthm}

%% The lineno packages adds line numbers. Start line numbering with
%% \begin{linenumbers}, end it with \end{linenumbers}. Or switch it on
%% for the whole article with \linenumbers.
%% \usepackage{lineno}

\journal{Robotics and Autonomous Systems}

\begin{document}

\begin{frontmatter}

\title{Multimodal robot with bio-inspired gait arbitration across simulated and real environments} %% Article title

\author{Samuel Federico Carlos-Rodríguez\fnref{depA}}
\author{Diego Alejandro Caro-Vaca\fnref{depA}}
\author{David Caicedo-Samboní\fnref{depA}}
\author{Sebastián Grajales-Pulgarín\fnref{depA}}
\author{Julián Hurtado-López\corref{cor1}}
\ead{jhurtado@uao.edu.co}
\cortext[cor1]{Corresponding author}
\author{David Fernando Ramírez-Moreno}

\affiliation{organization={Faculty of Engineering and Basic Sciences, Universidad Autónoma de Occidente},
            addressline={Cll 25 \# 115-85 Km 2 Vía Cali - Jamundí},
            city={Cali},
            postcode={760030},
            state={Valle},
            country={Colombia}}

\fntext[depA]{Mechatronic Engineering}

%% Abstract
\begin{abstract}
This work presents the design, simulation, and implementation of a bio-inspired quadruped robotic platform equipped with a hybrid wheel-leg locomotion system and controlled through a neural architecture modeled after basal ganglia mechanisms. The platform integrates sensory modules for environmental perception, arbitration modules for behavioral decision-making, and locomotor modules for gait generation, enabling adaptive responses to diverse stimuli and terrain conditions. Simulation and implementation demonstrate the robot's ability to switch autonomously between locomotion modes ---quadrupedal walking, differential rolling, and omnidirectional motion--- when facing obstacles, appetitive stimuli, or aversive conditions. \revblue{In simulation, an ablation study showed that removing the circuit's lateral inhibition or its STN/STR layers reduced success under stimulus conflict from 100\% to 25\% and 33\%, respectively. Against a priority-rule finite state machine (FSM) baseline under matched conditions, the neural arbitration succeeded in every three-stimulus conflict trial, where the FSM succeeded in none, and reached 85\% success versus 45\% in an inspection task with an initially hidden target.} Across simulated and real-robot trials, the controller consistently avoided critical instability and produced terrain-specific, repeatable locomotion signatures, laying the groundwork for scaling toward robust real-world infrastructure inspection and operation in complex, unstructured environments.
\end{abstract}

%%Graphical abstract
\begin{graphicalabstract}
\includegraphics[width=\textwidth]{GRAPHICAL_ABSTRACT}
\end{graphicalabstract}

%%Research highlights
\begin{highlights}
\item \revblue{Hybrid wheel-leg quadruped with three gait modes, tested in simulation and hardware}
\item \revblue{A basal-ganglia winner-take-all circuit arbitrates gaits from camera, LiDAR and IMU}
\item \revblue{Ablating lateral inhibition or STN/STR cuts simulated conflict success to 25--33\%}
\item \revblue{In three-stimulus conflict, neural arbitration succeeds in 100\% of trials; FSM in 0\%}
\item \revblue{Inspection task with an initially hidden target: 85\% success vs.\ 45\% for the FSM}
\end{highlights}

%% Keywords
\begin{keyword}

Wheel-legged quadruped \sep Hybrid locomotion \sep Gait mode switching \sep Arbitration \sep Basal ganglia model \sep Multimodal sensing

\end{keyword}

\end{frontmatter}

%% Add \usepackage{lineno} before \begin{document} and uncomment 
%% following line to enable line numbers
%% \linenumbers

%% main text
%%
